% Generated draft: requires mathematical and source review.
\documentclass[11pt]{article}
\usepackage[margin=25mm]{geometry}
\usepackage{fontspec}
\setmainfont{Latin Modern Roman}
\usepackage{amsmath,amssymb,mathtools,unicode-math}
\setmathfont{Latin Modern Math}
\usepackage{enumitem,needspace,xurl,hyperref}
\setlength{\parindent}{0pt}
\setlength{\parskip}{5pt}
\setlength{\emergencystretch}{3em}
\newcommand{\N}{\mathbb{N}}
\newcommand{\Z}{\mathbb{Z}}
\newcommand{\Q}{\mathbb{Q}}
\newcommand{\R}{\mathbb{R}}
\newcommand{\C}{\mathbb{C}}
\newcommand{\F}{\mathbb{F}}
\newcommand{\A}{\mathbb{A}}
\newcommand{\PP}{\mathbb{P}}
\newcommand{\E}{\mathbb{E}}
\newcommand{\eps}{\varepsilon}
\newcommand{\dd}{\,\mathrm{d}}
\begin{document}
\section*{UnsolvedMath ID: 30002353}
\textbf{Source problem code:} catalogue entry
\subsection{Definitions and mathematical statement}
Let \(k\) be a number field, that is, a finite extension of \(\Q\), and fix an algebraic closure \(\bar k\). Write \(G_k=\operatorname{Gal}(\bar k/k)\). Use the standard convention that a hyperbolic curve \(X/k\) is a smooth, geometrically connected curve whose smooth projective compactification \(\overline X\) has genus \(g\) and whose geometric boundary \((\overline X\setminus X)(\bar k)\) has cardinality \(r\), with \(2g-2+r>0\). This convention allows both proper and nonproper curves; the supplied statement imposes no further restriction.

Choose a geometric base point \(\bar x\) of \(X\). The arithmetic étale fundamental group fits into the exact sequence of profinite groups
\[
1\longrightarrow \pi_1^{\mathrm{\acute et}}(X_{\bar k},\bar x)
\longrightarrow \pi_1^{\mathrm{\acute et}}(X,\bar x)
\xrightarrow{p}G_k\longrightarrow 1.
\]
A Galois section is a continuous group homomorphism
\[
s:G_k\longrightarrow\pi_1^{\mathrm{\acute et}}(X,\bar x),
\qquad p\circ s=\operatorname{id}_{G_k}.
\]
Sections are considered up to conjugation by \(\pi_1^{\mathrm{\acute et}}(X_{\bar k},\bar x)\).

Write \(K=k(X)\) for the function field. A valuation of \(K\) is a map \(w:K\to\Gamma\cup\{\infty\}\), where \(\Gamma\) is a totally ordered abelian group, such that \(w(0)=\infty\), \(w(a)\in\Gamma\) for \(a\ne0\), \(w(ab)=w(a)+w(b)\) for nonzero \(a,b\), and \(w(a+b)\ge\min\{w(a),w(b)\}\).

\textbf{Source question (valuative formulation attributed to Grothendieck's section conjecture).} For every number field \(k\), every hyperbolic curve \(X/k\), and every continuous Galois section \(s\) of the displayed exact sequence, does there exist a suitable \(k\)-valuation \(w\) of \(k(X)\) that induces the conjugacy class \([s]\)?

\textbf{Unresolved qualification.} The supplied record does not define ``suitable \(k\)-valuation'' or the relation by which such a valuation induces a section of this étale fundamental group sequence. In particular, it does not specify the restriction of \(w\) to \(k\), its rank, residue-field conditions, allowed centers, or the required relation with decomposition groups. The general valuation definition above supplies notation only; it does not determine the admissible valuations. Consequently, the question retains these source terms explicitly and is not yet a fully specified valuative assertion. No equivalence with a rational-point formulation, and no replacement by a birational or local statement, is asserted here.
\subsection{Short English statement}
The supplied question asks whether every continuous Galois section of a hyperbolic curve over a number field comes, up to geometric conjugacy, from a suitable valuation of its function field. The source leaves both suitability and the meaning of ``comes from'' unspecified, so the valuative target requires clarification. Its open status is a supplied curation label.
\subsection{Sources}
\begin{itemize}
\item Supplied local UnsolvedMath registry record, permanent ID \(30002353\), title \emph{Valuative Form of Grothendieck's Section Conjecture}, exact source problem code \texttt{catalogue entry}, in the \emph{Oberwolfach Reports Open Problems} collection. The record attributes the conjecture to Grothendieck, lists proposed year \(2013\), and names no individual proposer.
\item Supplied report citation: \emph{The Arithmetic of Fields (2014)}, Oberwolfach Reports, DOI \texttt{10.4171/owr/2013/30}, pp.\ 1771--1818. \url{https://doi.org/10.4171/owr/2013/30}. The supplied title year is retained alongside the DOI's \(2013\) component; their relationship has not been independently checked. No problem-specific page is supplied.
\item Stix, \emph{On Galois sections for hyperbolic p-adic curves}. \url{https://www.math.uni-frankfurt.de/~stix/research/preprints/owr-STIX-Jun13.pdf}.
\item Pop--Stix, \emph{Arithmetic in the fundamental group of a p-adic curve}. \url{https://arxiv.org/abs/1111.1354}.
\item Saidi, \emph{Section conjecture over function fields}. \url{https://arxiv.org/abs/1512.01207}.
\item Supplied literature assessment, checked at the exact recorded date \texttt{2026-08-21}: the number-field Grothendieck section conjecture remains unresolved; valuative and p-adic analogues are proved; conditional reductions over finitely generated fields do not establish the number-field assertion. This is supplied dated curation derived from the OWR record, not an independently verified assessment. The record itself requests independent verification of time-sensitive status and cited claims; the scope and hypotheses of the claimed analogues are not specified here.
\item Supplied verification metadata labels the statement \texttt{corrected\_verified} and says it was checked against the official Oberwolfach report. This is a source-provided verification claim, not verification performed for this draft. The supplied problem status is \texttt{open}.
\end{itemize}
\end{document}
